% sections/12_api_contract.tex
\section{API Contract and Serving Guarantees}\label{sec:api-contract}

This section specifies the programmatic interface through which consumers
access published \ISI data.  The contract defines endpoint semantics,
response formats, caching behaviour, error taxonomy, and backward-compatibility
rules.  The API is a \textbf{read-only} serving layer; it does not expose
computation, ingestion, or administrative functions.

\subsection{Transport and Protocol}\label{sec:api:transport}

\begin{itemize}
  \item \textbf{Protocol.}  HTTPS (TLS~1.2 or later).  Plaintext HTTP
    requests receive a \texttt{301 Moved Permanently} redirect to the HTTPS
    endpoint.
  \item \textbf{Content type.}  All responses are served as
    JSON (\texttt{application/json}) with UTF-8 encoding unless
    otherwise specified per endpoint.
  \item \textbf{Authentication.}  The public API is unauthenticated.  Rate
    limiting is enforced per source IP address (see
    \cref{sec:api:rate-limiting}).
  \item \textbf{Base URL.}  The canonical base URL is published in the
    project repository.  All endpoint paths in this section are relative to
    the base URL.
\end{itemize}

\subsection{Endpoint Taxonomy}\label{sec:api:endpoints}

The API exposes the following endpoints.  All endpoints accept only
\texttt{GET} requests.  Other HTTP methods receive
\texttt{405 Method Not Allowed}.

\subsubsection{\texttt{GET /v1/snapshots}}

Returns the list of available snapshots.

\begin{table}[htbp]
  \centering\small
  \caption{Response schema for \texttt{GET /v1/snapshots}.}%
  \label{tab:api-snapshots}
  \begin{tabularx}{\textwidth}{@{}l l X@{}}
    \toprule
    \textbf{Field} & \textbf{Type} & \textbf{Description} \\
    \midrule
    \texttt{snapshots}   & \texttt{obj[]} & Array of snapshot summary objects. \\
    \texttt{.vintage\_year} & \texttt{int} & Vintage year. \\
    \texttt{.methodology\_version} & \texttt{str} & Version string. \\
    \texttt{.published\_at} & \texttt{str} & ISO~8601 publication timestamp. \\
    \texttt{.country\_count} & \texttt{int} & Number of countries in the snapshot. \\
    \bottomrule
  \end{tabularx}
\end{table}

\subsubsection{\texttt{GET /v1/snapshots/\{vintage\}/scores}}

Returns country-level scores for the specified vintage.

\paragraph{Path parameters.}
\texttt{\{vintage\}}: four-digit vintage year (e.g., \texttt{2024}).

\paragraph{Query parameters.}
\begin{description}[style=nextline, leftmargin=5em, labelwidth=4.5em]
  \item[\texttt{country}] Optional.  Comma-separated list of ISO~3166-1
    alpha-2 codes.  Filters the response to the specified countries.
  \item[\texttt{format}]  Optional.  One of \texttt{json} (default) or
    \texttt{csv}.  When \texttt{csv}, the response content type is
    \texttt{text/csv;\,charset=utf-8} and the body contains the
    \texttt{country\_scores.csv} content directly.
\end{description}

\paragraph{Response.}
The JSON response body is an object with a \texttt{scores} array whose
elements follow the column schema of \texttt{country\_scores.csv}
(\cref{tab:schema-country-scores}).  All numeric scores are returned as
strings (fixed-point, 8~decimal places).

\subsubsection{\texttt{GET /v1/snapshots/\{vintage\}/axes}}

Returns axis-level scores for all countries in the specified vintage.
The response schema mirrors \texttt{axis\_scores.csv}
(\cref{tab:schema-axis-scores}).

\subsubsection{\texttt{GET /v1/snapshots/\{vintage\}/channels}}

Returns channel-level scores.  The response schema mirrors
\texttt{channel\_scores.csv} (\cref{tab:schema-channel-scores}).

\subsubsection{\texttt{GET /v1/snapshots/\{vintage\}/integrity}}

Returns the integrity artefacts: country-level hashes, the snapshot hash,
the manifest hash, and the Ed25519 signature.  This endpoint provides
sufficient information for a consumer to perform the full verification
procedure (\cref{sec:integrity:verify}) without downloading the snapshot
archive.

\subsubsection{\texttt{GET /v1/snapshots/\{vintage\}/registry}}

Returns the methodology registry (\cref{sec:snapshot:registry}) as a
JSON object.

\subsubsection{\texttt{GET /v1/scenario}}

Executes a scenario simulation and returns the results.

\paragraph{Query parameters.}
\begin{description}[style=nextline, leftmargin=5em, labelwidth=4.5em]
  \item[\texttt{vintage}] Required.  Four-digit vintage year.
  \item[\texttt{country}] Required.  Single ISO~3166-1 alpha-2 code.
  \item[\texttt{a1} \ldots\ \texttt{a6}] Optional.  Axis-level adjustment
    factors, each a decimal in \([-0.20, +0.20]\).  Omitted axes default
    to~\(0\).
\end{description}

\paragraph{Response.}
The response contains the baseline scores, simulated scores, baseline rank,
simulated rank, and the applied adjustment vector.  All scores are returned
as fixed-point strings.

\subsection{ETag and Conditional Requests}\label{sec:api:etag}

Every response includes an \texttt{ETag} header whose value is the
hex-encoded SHA-256 hash of the response body.  Clients may issue
conditional requests using \texttt{If-None-Match}; the server returns
\texttt{304 Not Modified} if the ETag matches.

Because snapshot data is immutable (\cref{sec:snapshot:immutability}),
ETags for snapshot endpoints are stable: the same request to the same
vintage always produces the same ETag.  This property is a direct
consequence of determinism invariant \textbf{D-1}
(\cref{sec:det:invariants}).

Responses also include a \texttt{Cache-Control} header:
\begin{itemize}
  \item Snapshot endpoints: \texttt{Cache-Control:\ public,\ max-age=86400,\ immutable}.
  \item Scenario endpoint: \texttt{Cache-Control:\ public,\ max-age=3600}.
  \item Snapshot-list endpoint: \texttt{Cache-Control:\ public,\ max-age=3600}.
\end{itemize}

The \texttt{immutable} directive on snapshot endpoints signals to
intermediary caches that the response will never change for the given URL,
reflecting the immutability contract (\cref{sec:snapshot:immutability}).

\subsection{Error Taxonomy}\label{sec:api:errors}

The API uses the following error response structure:

\begin{quote}
\ttfamily\small
\{"error": \{"code": "<CODE>", "message": "<human-readable>"\}\}
\end{quote}

\Cref{tab:api-errors} enumerates the error codes.

\begin{longtable}{@{}l l l >{\raggedright\arraybackslash}p{5.5cm}@{}}
  \caption{API error codes.}\label{tab:api-errors} \\
  \toprule
  \textbf{HTTP} & \textbf{Code} & \textbf{Scope} & \textbf{Condition} \\
  \midrule
  \endfirsthead
  \caption[]{API error codes (continued).} \\
  \toprule
  \textbf{HTTP} & \textbf{Code} & \textbf{Scope} & \textbf{Condition} \\
  \midrule
  \endhead
  \midrule
  \multicolumn{4}{r@{}}{\footnotesize\itshape Continued on next page} \\
  \endfoot
  \bottomrule
  \endlastfoot
  400 & \texttt{INVALID\_VINTAGE}     & Path     & Vintage year is not a
    four-digit integer or is outside the range of published snapshots. \\
  400 & \texttt{INVALID\_COUNTRY}     & Query    & Country code is not a
    valid ISO~3166-1 alpha-2 code. \\
  400 & \texttt{INVALID\_ADJUSTMENT}  & Query    & Scenario adjustment factor
    is outside \([-0.20, +0.20]\) or is not a valid decimal. \\
  400 & \texttt{MISSING\_PARAMETER}   & Query    & A required query parameter
    is absent. \\
  404 & \texttt{VINTAGE\_NOT\_FOUND}  & Path     & The requested vintage year
    does not correspond to a published snapshot. \\
  404 & \texttt{COUNTRY\_NOT\_FOUND}  & Query    & The requested country code
    is not present in the specified vintage. \\
  405 & \texttt{METHOD\_NOT\_ALLOWED} & Method   & HTTP method other than
    \texttt{GET} was used. \\
  429 & \texttt{RATE\_LIMITED}         & Client   & Client has exceeded the
    rate limit. \\
  500 & \texttt{INTERNAL\_ERROR}      & Server   & An unexpected server-side
    error occurred.  The response includes a correlation identifier for
    incident tracking. \\
\end{longtable}

\subsection{Rate Limiting}\label{sec:api:rate-limiting}

Rate limiting is enforced per source IP address using a token-bucket
algorithm.  The default limits are:

\begin{table}[htbp]
  \centering\small
  \caption{Default rate limits.}\label{tab:api-rate-limits}
  \begin{tabular}{@{}l l l@{}}
    \toprule
    \textbf{Endpoint group} & \textbf{Requests/minute} & \textbf{Burst} \\
    \midrule
    Snapshot data endpoints   & 60  & 10 \\
    Scenario endpoint         & 30  & 5 \\
    Snapshot list              & 60  & 10 \\
    \bottomrule
  \end{tabular}
\end{table}

When a client exceeds the rate limit, the server returns
\texttt{429 Too Many Requests} with a \texttt{Retry-After} header
indicating the number of seconds until the next permitted request.

\subsection{Versioned Endpoint Prefix}\label{sec:api:versioning}

All endpoints are prefixed with a version identifier (\texttt{/v1/}).  The
API version is \emph{independent} of the methodology version: API~v1 may
serve snapshots from methodology versions~1.0, 1.1, 1.2, etc.

A new API major version (\texttt{/v2/}) is introduced only when the response
schema undergoes a backward-incompatible change.  When a new API version is
introduced, the previous version remains operational for a minimum of
twelve months with a \texttt{Sunset} header indicating the deprecation date.

\subsection{Backward-Compatibility Rules}\label{sec:api:compat}

Within a given API major version, the following changes are
backward-compatible and do not require a version increment:

\begin{itemize}
  \item Adding new optional query parameters.
  \item Adding new fields to JSON response objects.
  \item Adding new snapshot vintages.
  \item Adding new warning codes to the \texttt{warnings} field.
\end{itemize}

The following changes are \textbf{breaking} and require a new API major
version:

\begin{itemize}
  \item Removing or renaming an existing response field.
  \item Changing the type or semantics of an existing response field.
  \item Removing an endpoint.
  \item Changing the serialisation format of scores (e.g., from string to
    number).
  \item Altering the ETag computation algorithm.
\end{itemize}

\subsection{Content Negotiation}\label{sec:api:content}

Clients may request CSV output via the \texttt{format=csv} query parameter
on score endpoints.  The CSV output is byte-identical to the corresponding
file in the snapshot archive, ensuring that consumers who download the
snapshot and consumers who use the API receive the same data.

JSON is the default format.  The \texttt{Accept} header is not used for
content negotiation; the \texttt{format} query parameter takes precedence.

\subsection{Health and Readiness}\label{sec:api:health}

The API exposes two operational endpoints that are excluded from the
versioned prefix:

\begin{description}[style=nextline]
  \item[\texttt{GET /healthz}]
    Returns \texttt{200~OK} with body \texttt{\{"status":"ok"\}} if the
    server process is running and able to accept connections.  This endpoint
    performs no database or file-system checks.  It is intended for load
    balancers and container orchestrators.
  \item[\texttt{GET /readyz}]
    Returns \texttt{200~OK} with body \texttt{\{"status":"ready"\}} if
    the server has loaded at least one snapshot into memory and is able to
    serve score requests.  If no snapshot is loaded, it returns
    \texttt{503~Service Unavailable} with
    \texttt{\{"status":"not\_ready"\}}.
\end{description}

Neither endpoint is subject to rate limiting or authentication.  Neither
endpoint returns cache headers.  Both endpoints are excluded from the
API versioning scheme and may not be removed or renamed.

\subsection{Cross-Origin Resource Sharing}\label{sec:api:cors}

The API sets the following CORS headers on all responses:

\begin{itemize}
  \item \texttt{Access-Control-Allow-Origin:~*}
  \item \texttt{Access-Control-Allow-Methods:~GET,~OPTIONS}
  \item \texttt{Access-Control-Max-Age:~86400}
\end{itemize}

This permits browser-based clients to access the API directly without a
server-side proxy.
